\centerline{\Huge \bf  Тренировочная Олимпиада 4}
\centerline{\Huge \bf 7 класс}
\centerline{\Huge \bf Уровень: Заключительный отбор в Сириус.}


\begin{flushright}
    \scriptsize{Глава 4. Четвёртое дитя.}
\end{flushright}

\problem{1}
Дата - 7 ноября 2016г.\\
Время - Утро 5:00.\\
Место - ОЦ Сириус.\\
Проблема - Пока АА спит, Расул Владимирович проснулся и ему нечего делать до занятий, которые начинаются в 10:00.\\
Решение - Сразу после пробуждения Расул Владимирович начал записывать в тетрадку округленные в большую сторону целые значения минут, когда часовая стрелка и минутная лежали на одной прямой, но были противоположно направленны.

Найдите сумму значений, записанных в тетрадку Расулом Владимировичем до начала занятий в Сириусе.
\begin{flushright}
    \textit{\small{(Доржеев А.А.)}}
\end{flushright}

\problem{2}
Каждый из 14 учеников Расула Владимировича кинул приглашение в пати ровно $k$ другим ученикам. Если какие-то два ученика кинули приглашение друг другу, то они попадают в пати и идут вместе играть. При каком наименьшем $k$ какие-то два ученика пойдут играть?
\begin{flushright} 
    \textit{\small{(Гасанов Р.В.)}}
\end{flushright}

\problem{3}
В бесконечной аллее в ряд выстроены пронумерованные (начиная с $1$ и по порядку) деревья на расстоянии $1$ метр друг от друга. Заира прогуливалась по этой аллее и тянула за собой километровый шарф. Докажите, что в какой-то момент её шарф окутает ровно $10$ деревьев с простыми номерами.

\begin{flushright}
    \textit{\small{(Гасанов Р.В.)}}
\end{flushright}

\problem{4}
Биссектриса внешнего угла $C$ треугольника $ABC$ пересекает биссектрису внутреннего угла $B$ в точке $K$. Оказалось, что $CB = CK$. Докажите, что треугольник $ABC$ равнобедренный.
\begin{flushright}
    \textit{\small{(Гасанов Р.В.)}}
\end{flushright}

\newpage

\hspace{3mm}

\problem{5}
Про положительные числа $a, b, c$ известно, что $abc = 1$ и $a^3 \geq 12$. Докажите, что
\[\dfrac{a^2}{2} + b^2 + c^2 \geq ab + bc + ac.\]

\begin{flushright}
    \textit{\small{(Гасанов Р.В.)}}
\end{flushright}

\problem{6}
У красноухого бабуина Вячеслава есть $N$ книг Кормана по алгоритмам. Однажды Вячеслав забрался на 36-этажное главное здание МГУ и начал скидывать оттуда книги Кормана с разных этаже. Если книга не разрывалась при приземлении, то он подбирал её и мог скинуть ещё раз с какого-нибудь этажа ГЗ МГУ, а если всё же разорвалась, то увы скидывать он её больше не сможет. За какое минимальное количество бросков бабуин сможет наверняка понять, начиная с какого этажа книга по алгоритмам полностью разрывается?
\begin{flushright}
    \textit{\small{(Гасанов Р.В.)}}
\end{flushright}